\section{Introduction}
\label{sec:intro}

Bayesian inverse problems are most naturally formulated on function
spaces.  A prior is placed on an unknown $u$ in a separable Hilbert
space $\HH$, data are generated through a forward map, and Bayes'
formula defines a posterior measure rather than a finite-dimensional
density.  This viewpoint makes stability under perturbations and
discretization part of the mathematical problem, not merely part of an
implementation \citep{stuart2010,dashti2017}.

Learned generative models enlarge the class of priors used in inverse
problems.  Score-based constructions are particularly attractive
because they combine a Gaussian reference process with a learned
vector field.  In function spaces, however, a learned vector field is
not automatically an admissible stochastic drift, and a drift used in
a generative SDE is not automatically the score of the SDE's own
marginal law.  Reverse-time score equations have a specific sign,
time orientation, and conditional-expectation structure
\citep{pidstrigach2023}.  Confusing these objects can invalidate an
otherwise plausible posterior argument.

This paper studies a narrower object that can be defined without that
identification.  Starting from the stationary Gaussian
Ornstein--Uhlenbeck process with covariance $\mathcal C$, we add a
measurable drift $\drift(t,u)$ taking values in the Cameron--Martin
space of $\mathcal N(0,\mathcal C)$.  The terminal law is the prior.
The parameterization may be learned using score-matching data, but the
mathematical assumptions concern the resulting drift itself.  A
finite-horizon Girsanov condition then relates the prior to the
reference Gaussian.

The first result is posterior well-posedness.  For bounded linear
observations and fixed positive noise level, existence and local
Hellinger stability require only a finite second moment of the prior.
Girsanov is therefore not part of Bayes' formula; it is a separate
device for proving properties of this particular prior construction.
This separation keeps the statistical statement valid even when a
different method is used to construct the prior.

The second result concerns spectral approximation.  If a
finite-dimensional approximation is embedded in $\HH$ by setting all
unresolved coordinates to zero, its law is singular with respect to a
nondegenerate Gaussian measure on $\HH$.  Consequently, its Hellinger
distance from the full prior is maximal, regardless of how small the
discarded covariance eigenvalues are.  We instead use a
\emph{full-space lift}: the learned drift and the likelihood are
projected onto the first $N$ modes, while the unresolved modes retain
their reference Gaussian law.  This produces a genuine
infinite-dimensional probability measure and yields the bound
\[
  \dH(\muyN,\muy)
  \leq C_R\!\left(\sqrt{\varepsilon_N}
                  +\Bigl(\sum_{j>N}\lambda_j\Bigr)^{1/2}\right).
\]
The first term is a Cameron--Martin energy error between path drifts;
the second is the root trace of the unresolved covariance.  For
$\lambda_j\asymp j^{-2s}$ with $s>1/2$, the latter is of order
$N^{1/2-s}$, not $\sqrt{\lambda_{N+1}}$.

Posterior contraction needs more information than posterior
well-posedness.  It depends on the observation experiment, the
regularity of the truth, prior mass near that truth, tests, and an
identifiable loss.  Even in the conjugate Gaussian sequence model, the
rate depends on the prior regularity, truth regularity, degree of
ill-posedness, and prior scaling \citep{knapik2011}.  We therefore give
a conditional transfer criterion rather than assigning a universal
rate to the learned prior.  A local lower bound on the Girsanov density
transfers reference-Gaussian mass to the learned prior, but a global
$L^2$ density bound alone does not provide that local information.

For every fixed noise level, the posterior is locally Hellinger
Lipschitz in the observations.  The lifted approximation converges when
both the drift energy and covariance tail vanish.  Statistical
contraction then requires separate local-mass, sieve, and testing
arguments for the observation model.  The identifiable parameter is
the equivalence class induced by the forward map; stronger recovery
requires injectivity or a quantitative inverse stability estimate on
the relevant parameter class.

These statements complement work devoted to constructing or analysing
diffusion-based samplers in function spaces
\citep{baldassari2025langevin,baker2026,cheng2026}.  They do not provide
a sampler convergence theorem, a nonlinear-forward-map theorem, or an
unconditional contraction rate.  Those questions require assumptions
not supplied by the present model.

\paragraph{Organization.}
\Cref{sec:setup} defines the prior, likelihood, and Hellinger metric.
\Cref{sec:wellposed} proves posterior stability and states the lifted
discretization theorem.  \Cref{sec:contraction} isolates the additional
conditions needed for contraction and identifiability.
\Cref{sec:experiments} records the limited role of the accompanying
code, and \Cref{sec:discussion} summarizes the remaining mathematical
gaps.  Detailed proofs are in the appendices.
